\chapter*{Sommario} % senza numerazione
\label{sommario}

\addcontentsline{toc}{chapter}{Sommario} % da aggiungere comunque all'indice

L'interazione tra linguaggio naturale e database relazionali rappresenta
una delle sfide più rilevanti nell'ambito dell'Intelligenza Artificiale
applicata. Questa tesi presenta lo sviluppo di un framework avanzato di
Text-to-SQL progettato per ottimizzare la generazione di query
attraverso un approccio modulare e iterativo.

A differenza dei modelli tradizionali, il sistema implementato si basa
su un'architettura a doppia orchestrazione: lo
\texttt{SchemaOrchestrator}, dedicato alla ricostruzione semantica e
all'acquisizione della struttura del database, e il
\texttt{QueryOrchestrator}, responsabile della generazione e
raffinamento del codice SQL. Un elemento distintivo della ricerca è
l'integrazione di tecniche di Retrieval-Augmented Generation (RAG)
tramite componenti di storage vettoriale (\texttt{SchemaStore} e
\texttt{QueryStore}), che permettono al sistema di mantenere una memoria
storica delle interrogazioni passate e dei relativi fallimenti.

Il cuore innovativo del lavoro risiede nel ciclo di feedback
automatizzato: il sistema non si limita a produrre un output, ma esegue
una validazione sintattica e semantica in tempo reale. In caso di
errore, un modulo di analisi (\texttt{LLMFeedback}) categorizza
l'anomalia, fornendo istruzioni mirate al modello per la correzione
autonoma. I risultati dimostrano come l'uso combinato
di prompt ingegnerizzati, contestualizzazione dinamica dello schema e
apprendimento dagli errori pregressi aumenti significativamente
l'affidabilità delle query generate, portando in media il \(21,54\%\) di successi in più, riducendo la distanza tra l'intento
dell'utente e l'esecuzione tecnica su database eterogenei.

\section*{Uso di strumenti di Intelligenza Artificiale}

Per lo sviluppo del progetto di tesi sono stati impiegati strumenti di Intelligenza Artificiale,
in particolare durante la fase di elaborazione del codice è stato utilizzato ChatGPT \cite{chatgpt} ed il suo tool integrato
Codex \cite{openai_codex} per organizzare la struttura e la generazione di alcune formule ausiliarie. 
Per quanto riguarda la stesura del testo, è stato sfruttato NotebookLLM \cite{google_notebooklm} per la raccolta, gestione e traduzione delle fonti
mentre l'utilizzo di Google Gemini \cite{google_gemini} ha permesso di controllare la correttezza dal punto di vista grammaticale e ortografico del testo.
